\documentclass[11pt]{article}
\usepackage[paperwidth=11in,paperheight=8.5in,margin=0.25in]{geometry}
\usepackage[T1]{fontenc}
\usepackage[utf8]{inputenc}
\usepackage[english]{babel}
\usepackage{lmodern}
\usepackage{microtype}
\usepackage{amsmath,amssymb}
\usepackage{xcolor}
\usepackage{tikz}
\usetikzlibrary{calc}
\hyphenpenalty=8000
\exhyphenpenalty=8000
\emergencystretch=2em
\definecolor{QCBlue}{HTML}{1F4E79}
\definecolor{QCGreen}{HTML}{2E6B5F}
\definecolor{QCOrange}{HTML}{B65B2A}
\definecolor{QCPurple}{HTML}{4B3F72}
\definecolor{QCLight}{HTML}{F4F6F8}
\definecolor{QCBorder}{HTML}{D0D6DC}
\definecolor{QCText}{HTML}{111111}
\newlength{\Margin} \setlength{\Margin}{0.25in}
\newlength{\Gap} \setlength{\Gap}{0.1in}
\newlength{\TitleH} \setlength{\TitleH}{0.8in}
\newlength{\HeaderH} \setlength{\HeaderH}{0.38in}
\newlength{\Pad} \setlength{\Pad}{0.1in}
\newlength{\UsableW}
\newlength{\CellW}
\newlength{\CellH}
\setlength{\UsableW}{\dimexpr\paperwidth-2\Margin\relax}
\setlength{\CellW}{5.2in}
\setlength{\CellH}{3.5in}
\newcommand{\QuadBox}[4]{%
  \node[anchor=north west, draw=QCBorder, line width=0.6pt, rounded corners=10pt, fill=white, minimum width=\CellW, minimum height=\CellH, inner sep=0pt] (B) at #1 {};
  \fill[#2, rounded corners=10pt] ([xshift=0pt,yshift=0pt]B.north west) rectangle ([xshift=0pt,yshift=-\HeaderH]B.north east);
  \node[anchor=west, text=white, font=\bfseries] at ([xshift=0.18in,yshift=-0.22in]B.north west) {#3};
  \node[anchor=north west, text=QCText] at ([xshift=\Pad,yshift=-\HeaderH-\Pad]B.north west) {\begin{minipage}[t]{\dimexpr\CellW-2\Pad\relax}\raggedright\small #4\end{minipage}};
}
\pagestyle{empty}
\begin{document}
\begin{tikzpicture}[remember picture,overlay]
\coordinate (NW) at ([xshift=\Margin,yshift=-\Margin]current page.north west);
\node[anchor=north west, draw=QCBorder, line width=0.6pt, rounded corners=12pt, fill=QCLight, minimum width=\UsableW, minimum height=\TitleH, inner sep=0pt] (T) at (NW) {};
\node[anchor=north west] at ([xshift=0.22in,yshift=-0.22in]T.north west) {\begin{minipage}[t]{\dimexpr\UsableW-0.44in\relax}{\Huge\bfseries \textcolor{QCBlue}{Method of Images Formula Sheet}}\par\vspace{2pt}{\normalsize \textbf{Diego Juarez}\hfill \textbf{Computational Electrodynamics}\hfill \textbf{\today}}\end{minipage}};
\coordinate (GridNW) at ([yshift=-0.15in]T.south west);
\coordinate (TL) at (GridNW);
\coordinate (TR) at ([xshift=\CellW+\Gap]GridNW);
\coordinate (BL) at ([yshift=-\CellH-\Gap]GridNW);
\coordinate (BR) at ([xshift=\CellW+\Gap,yshift=-\CellH-\Gap]GridNW);
\QuadBox{(TL)}{QCBlue}{1) Core Principle}{
Find any potential in the physical region that:

1. Solves Poisson or Laplace there.

2. Satisfies the conductor boundary conditions.

Then uniqueness guarantees it is the physical solution.

Image charges are mathematical devices outside the physical region.
}
\QuadBox{(TR)}{QCGreen}{2) Grounded Plane}{
Real charge at $(0,0,a)$, image charge $-q$ at $(0,0,-a)$.
\[
\Phi(\mathbf r)=\frac{q}{4\pi\varepsilon_0}\left(
\frac{1}{|\mathbf r-\mathbf r_0|}-\frac{1}{|\mathbf r-\mathbf r_0^*|}
\right), \quad z>0
\]
Boundary:
\[
\Phi(x,y,0)=0
\]
Force on the real charge:
\[
F=\frac{1}{4\pi\varepsilon_0}\frac{q^2}{(2a)^2}
\]
Direction: toward the plane.
}
\QuadBox{(BL)}{QCOrange}{3) Grounded Sphere}{
Charge $q$ at distance $a>R$ from center.

Image charge:
\[
q'=-q\frac{R}{a}
\]
Image position:
\[
b=\frac{R^2}{a}
\]
Exterior potential:
\[
\Phi=\frac{1}{4\pi\varepsilon_0}\left(
\frac{q}{|\mathbf r-\mathbf r_0|}+\frac{q'}{|\mathbf r-\mathbf r_i|}
\right)
\]
with $\Phi(R,\theta)=0$.
}
\QuadBox{(BR)}{QCPurple}{4) Surface Charge and Checks}{
Induced surface charge on conductor:
\[
\sigma=\varepsilon_0 E_n
\]
or more carefully from the normal field just outside.

\textbf{Check list}

1. Image charges must lie outside the physical region.

2. Verify the boundary exactly, not approximately.

3. Use uniqueness before claiming the answer.

4. Grounded and isolated conductors are different problems.
}
\end{tikzpicture}
\end{document}
